\documentclass[11pt, a4paper]{report}

\usepackage[T1]{fontenc} 								
\usepackage[norsk]{babel}								
\usepackage[utf8]{inputenc}						
\usepackage{graphicx}       						
\usepackage{amsmath,amssymb}
\usepackage{grffile}
\usepackage{listings}
\usepackage{caption}
\usepackage[export]{adjustbox}
\usepackage{titling}
\setcounter{secnumdepth}{0}

\lstset{extendedchars=true, basicstyle=\footnotesize, numbers=left, numberstyle=\tiny, frame=shadowbox, tabsize=2, language=C, showstringspaces=false, breaklines=true,
  literate=
  {á}{{\'a}}1 {é}{{\'e}}1 {í}{{\'i}}1 {ó}{{\'o}}1 {ú}{{\'u}}1
  {Á}{{\'A}}1 {É}{{\'E}}1 {Í}{{\'I}}1 {Ó}{{\'O}}1 {Ú}{{\'U}}1
  {à}{{\`a}}1 {è}{{\`e}}1 {ì}{{\`i}}1 {ò}{{\`o}}1 {ù}{{\`u}}1
  {À}{{\`A}}1 {È}{{\'E}}1 {Ì}{{\`I}}1 {Ò}{{\`O}}1 {Ù}{{\`U}}1
  {ä}{{\"a}}1 {ë}{{\"e}}1 {ï}{{\"i}}1 {ö}{{\"o}}1 {ü}{{\"u}}1
  {Ä}{{\"A}}1 {Ë}{{\"E}}1 {Ï}{{\"I}}1 {Ö}{{\"O}}1 {Ü}{{\"U}}1
  {â}{{\^a}}1 {ê}{{\^e}}1 {î}{{\^i}}1 {ô}{{\^o}}1 {û}{{\^u}}1
  {Â}{{\^A}}1 {Ê}{{\^E}}1 {Î}{{\^I}}1 {Ô}{{\^O}}1 {Û}{{\^U}}1
  {ã}{{\~a}}1 {Ã}{{\~A}}1 {õ}{{\~o}}1 {Õ}{{\~O}}1
  {œ}{{\oe}}1 {Œ}{{\OE}}1 {æ}{{\ae}}1 {Æ}{{\AE}}1 {ß}{{\ss}}1
  {ç}{{\c c}}1 {Ç}{{\c C}}1 {ø}{{\o}}1 {å}{{\r a}}1 {Å}{{\r A}}1
  {€}{{\EUR}}1 {£}{{\pounds}}1}

\setlength{\textheight}{240mm} 
\setlength{\textwidth}{180mm}  
\topmargin -5mm 
\oddsidemargin -5mm
\captionsetup[figure]{singlelinecheck=off, justification = raggedleft} 
\pretitle{%
  \begin{center}
  \LARGE
  \includegraphics[width=6cm,height=6cm]{uitlogo.png}\\[\bigskipamount]
}
\posttitle{\end{center}}

\begin{document}

\title{Oblig 3 - Halvdupleksprotokoll}
\author{Fredrik Sandhei, Mathias Haukås \thanks{UiT - AUT-1001, obligatorisk innlevering 3}}
\date{\today}
\maketitle
\newpage
\tableofcontents
\newpage

\section{Introduksjon}
I samsvar med arbeidskravene til faget "Programmering med mikrokontrollere", trengs det tre arbeidskrav, i form av tre obligatoriske innleveringer, godkjent. I dette dokumentet ligger vår besvarelse på den siste obligatoriske innleveringen - oblig 3. 

\section{Hensikt}

\section{Teori}

\subsection{Protokoll}
Protokollen som skulle brukes i oppgaven er "halv-dupleks" - En datapakke sendes først fra datamaskinen til mikrokontrolleren ved hjelp RS-232 - kommunikasjon. Deretter forventer protokollen en transmit fra mikrokontrolleren. Dataen som sendes og mottas er forventet å være i ASCII-format. Hver Rx - prosedyre er kontrollert ved hjelp av rekkefølgen på pakkeelementene: Hver pakke begynner på en bokstav, enten $A$, $B$, $C$ eller $D$ etterfulgt med line feed, $\backslash n$. Når en Rx-prosedyre er gjennomført, avsluttes pakken med $R$ og $\backslash n$, og protokollen er klar for å motta en Tx-prosedyre, som ikke er like strengt lagt opp. Den forventer igjen det samme som i Rx, men rekkefølgen spiller ingen rolle, og ingen $R$ trengs ikke for å konkludere transmit. 

\subsection{USART}
For å opprette kommunikasjon mellom PC og displaykortet er USART - Universal Synchronous Asynchronous Receive Transmit - et nyttig verktøy. Det er en type seriell kommuikasjon, der dataelementene/bitsene shiftes over en kabel mellom enhetene. Ved hjelp av en ekstra kabel kan en bruke en seriell klokke mellom enhetene for å synkronisere dataregistreringen hos mottakeren. I vårt tilfelle benyttes ikke den serielle klokken. Kommunikasjonen blir dermed asynkron - UART. Asynkron USART bruker et start-bit og et (evt. to) stopp-bit i tillegg til den dataen som sendes for å bestemme overføringshastigheten. 

\section{Løsning}
\subsection{Implementasjon}

\section{Diskusjon}

\section{Konklusjon}
%\lstinputlisting{main.c}

\end{document}